%% LyX 2.3.6.1 created this file.  For more info, see http://www.lyx.org/.
%% Do not edit unless you really know what you are doing.
\documentclass[12pt,ngerman]{article}
\usepackage{mathptmx}
\usepackage[T1]{fontenc}
\usepackage[latin9]{inputenc}
\usepackage{geometry}
\geometry{verbose,tmargin=1cm,bmargin=2cm,lmargin=2cm,rmargin=2cm}
\usepackage{array}
\usepackage{textcomp}
\usepackage{amssymb}

\makeatletter

%%%%%%%%%%%%%%%%%%%%%%%%%%%%%% LyX specific LaTeX commands.
%% Because html converters don't know tabularnewline
\providecommand{\tabularnewline}{\\}

\@ifundefined{date}{}{\date{}}
%%%%%%%%%%%%%%%%%%%%%%%%%%%%%% User specified LaTeX commands.
\usepackage[colorlinks,
pdfpagelabels,
pdfstartview = FitH,
bookmarksopen = true,
bookmarksnumbered = true,
linkcolor = black,
plainpages = false,
hypertexnames = false,
citecolor = black] {hyperref}
\usepackage{tikz,pgf,pgfplots}
\usepackage{mathrsfs,float,multirow}
\usetikzlibrary{shapes.geometric,patterns,calc,matrix,arrows.meta}
\hypersetup{urlcolor=myred}
\setlength{\parindent}{0pt}
\usepackage{amsthm}
\usepackage{wrapfig}
\setlength\parindent{0pt}
\definecolor{myblue}{RGB}{0,56,101}
\definecolor{mygreen}{RGB}{106,176,35}
\definecolor{myred}{RGB}{176,47,44}
\definecolor{myred2}{RGB}{216,65,63}

%%%%% coloring tables %%%%
\definecolor{tablegreen}{RGB}{224,239,212}
\definecolor{tableyellow}{RGB}{255,251,204}
\definecolor{tablered}{RGB}{252,212,209}
\usepackage{colortbl,dcolumn}
\def\zz#1{%
\ifdim#1pt<0.5pt\cellcolor{tablered}\else
\ifdim#1pt<1pt\cellcolor{tableyellow}\else
\cellcolor{tablegreen}\fi\fi
#1}

%%%% tcolorbox %%%%%%%%
\usepackage[most]{tcolorbox}
\tcbuselibrary{breakable}
\setlength{\intextsep}{2pt}
\usepackage{pdfpages}

\newtcolorbox{kasten}[2]{
before skip=15pt,
after skip=15pt,
breakable,
enhanced jigsaw,
colback=myblue!5!white,
colframe=myblue,
colbacktitle=myblue,
opacityback=0.35,
attach boxed title to top left={yshift=-2mm},
fonttitle=\bfseries,
title=#1
}
\pagestyle{empty}

\newtcolorbox{mybox}[2][]{
enhanced,before skip=15pt,
after skip=15pt,
colback=black!5,
breakable,
colframe=black!50,
boxrule=0.2mm,
attach boxed title to top left={xshift=1cm,yshift*=1mm-\tcboxedtitleheight},
boxed title style={frame code={\path[fill=tcbcolback!30!black]([yshift=-1mm,xshift=-1mm]frame.north west)arc[start angle=0,end angle=180,radius=1mm]([yshift=-1mm,xshift=1mm]frame.north east)arc[start angle=180,end angle=0,radius=1mm];\path[left color=tcbcolback!60!black,right color=tcbcolback!60!black,middle color=tcbcolback!80!black]([xshift=-2mm]frame.north west) -- ([xshift=2mm]frame.north east)[rounded corners=1mm]-- ([xshift=1mm,yshift=-1mm]frame.north east)-- (frame.south east) -- (frame.south west)-- ([xshift=-1mm,yshift=-1mm]frame.north west)[sharp corners]-- cycle;},interior engine=empty,},
fonttitle=\bfseries,
title={#2},#1}

\numberwithin{table}{section}

\makeatother

\usepackage[font=small,labelfont=bf]{caption}
%%% Font
\usepackage{helvet}
\renewcommand{\familydefault}{\sfdefault}

\usepackage{babel}
\begin{document}
\pagenumbering{gobble}

\begin{tikzpicture}
	\draw[very thick,fill=myred,draw=myred,rounded corners] (-3.7,-.7) rectangle (8.8,.7);
	\filldraw[myred2] (-1,-.7)--(4,-.7)--(6,.7)--(1,.7)--cycle;
	\node[align=left,white] at (2,0) {\large Fachdidaktik 1 -- Grundlagen der Physikdidaktik};
	\filldraw[fill=white,draw=none] (10,-1.8) arc (250:110:2);
	\node at (11.4,0) {\includegraphics[scale=.07]{../../Uni-L-Logo.png}};
	
	\node[align=left] at (-1,-2) {
		\footnotesize{Prof.~Dr.~Philipp Bitzenbauer}\\
		\footnotesize{philipp.bitzenbauer@uni-leipzig.de}\\
		\footnotesize{Vor dem Hospitaltore 1, 1R13}\\
		\footnotesize{T +49 341 97-32772}
	};
	
	\node[align=left] at (5.5,-2) {
		\footnotesize{Dr.~Joaquin Veith}\\
		\footnotesize{joaquin.veith@uni-leipzig.de}\\
		\footnotesize{Vor dem Hospitaltore 1, 1R18}\\
		\footnotesize{T +49 341 97-32458}
	};
	%\draw (-3.7,-1.7)--(13.9,-1.7);
\end{tikzpicture}
\vspace{.5cm}
\begin{center}\bfseries\Large Informationsblatt zur Vortragsstruktur\end{center}
\vspace{1cm}
\small
Im Modul Fachdidaktik 1 werden ausgew�hlte Themen der Physikdidaktik mit besonderer Relevanz f�r die Unterrichtspraxis vertieft. Dazu halten alle Teilnehmenden einen Vortrag. Die Themen und Termine f�r die Vortr�ge werden in der ersten Sitzung am 13.10.2025 verteilt. Nachtr�gliche �nderungen sind nicht m�glich. Die ersten Vortr�ge finden am 10.11.2025 statt.\\
\\
Die Vortr�ge sind von ca. 20 min. Dauer und besitzen folgende Struktur:
\begin{itemize}
\item Theoretischer Hintergrund (ca. 10 min.): Aufbereitung des theoretischen Hintergrunds zum jeweiligen Thema, inklusive aktueller Forschungserkenntnisse.
\item Relevanz f�r die Unterrichtspraxis (ca. 10 min.): Herausarbeiten von Implikationen der theoretischen Fundierung und aktueller Forschungserkenntnisse f�r die Unterrichtspraxis.
\end{itemize}

Das grobe Einhalten der Vortragszeiten und -struktur ist Teil der Aufgabe. In ihren inhaltlichen Schwerpunkten und damit ihrem Aufbau unterscheiden sich die Vortr�ge je nach Thema. Die Referierenden besitzen die Freiheit, inhaltliche Schwerpunkte nach pers�nlichen Interessen zu setzen. Wichtig ist jedoch, dass zentrale Begriffe gekl�rt werden und wichtige theoretische Grundlagen vermittelt werden. Die Eigenleistung findet im zweiten Teil des Vortrags statt, wenn sie stichhaltig die Relevanz des Themas f�r den sp�teren Lehrberuf darlegen.\\
\\
Die Gestaltung des Vortrags beginnt daher stets mit einer ausf�hrlichen Literaturrecherche. Dazu bieten sich je nach Thema sowohl Lehrb�cher der Physikdidaktik (z.B. \href{https://doi.org/10.1007/978-3-642-41745-0}{Physikdidaktik. Theorie und Praxis von Kircher, Girwidz und H�u�ler, 2015} oder \href{https://www.friedrich-verlag.de/shop/physikdidaktik-kompakt-62784}{Physikdidaktik kompakt von Wiesner, Schecker und Hopf, 2014}) als auch Artikel aus wissenschaftlichen Zeitschriften an. Zum Start f�r Ihre Recherchen ist in der Zeitschiene im Moodle-Kurs, in der tabellarisch alle Themen und Daten des Moduls aufgeschl�sselt sind, zu jedem Thema ein besonders zentraler Fachartikel per Hyperlink hinterlegt. Empfehlenswert ist dar�ber hinaus ein Blick in wissenschaftliche Zeitschriften, insbesondere f�r den Einbezug aktueller Erkenntnisse. Eine Suche nach Artikeln in folgenden Zeitschriften kann besonders lohnenswert sein (Suchfunktion nutzen):
\begin{itemize}
	\item \href{https://www.tandfonline.com/journals/tsed20}{International Journal of Science Education}
	\item \href{https://journals.aps.org/prper/}{Physical Review Physics Education Research}
	\item \href{https://www.springer.com/journal/10763}{International Journal of Science and Mathematics Education}
	\item \href{https://iopscience.iop.org/journal/0031-9120}{Physics Education}
	\item \href{https://www.ejmste.com/}{Eurasia Journal of Mathematics, Science and Technology Education}
	\item \href{https://www.springer.com/journal/11165}{Research in Science Education}
\end{itemize}

Zur medialen Unterst�tzung der Vortr�ge gibt es keine Vorgaben. Ein guter Vortrag wird verstanden als eine Bereicherung f�r alle Teilnehmenden und f�r die Referierenden selbst:
\begin{itemize}
	\item Eine innovative Vortragsgestaltung,
	\item das Pr�sentieren eines theoretischen Hintergrunds zum Thema, der eine
	substanzielle Erweiterung der Vorlesung Fachdidaktik 1 darstellt und
	\item eine gut argumentierte Relevanz des Themas f�r den Physikunterricht
\end{itemize}
sind dabei entscheidender als das Abhaken formaler Vorgaben.\\
\\
Zur Pr�fungsvorleistung im Modul Fachdidaktik 1 geh�rt eine schriftliche Ausarbeitung des Vortrags. Diese entspricht einem Handout, das die Vortragsinhalte auf einer bis zwei DIN A4 Seiten zusammenfasst. Das Handout entspricht also einer �bersichtlichen und einheitlichen Zusammenfassung der wesentlichen Inhalte und soll einerseits f�r die Vorbereitung zur Modulpr�fung sowie andererseits f�r die sp�tere Unterrichtspraxis unterst�tzend wirken.\\
\\
Eine PDF der Vortragspr�sentation und des Handouts sind bis sp�testens zwei Tage vor Vortrag in den daf�r vorgesehenen Abgabemodulen im Moodlekurs hochzuladen. Jedem Vortrag folgt eine ca. 10 min�tige Diskussion im Plenum, in der offene Fragen und kritische Anmerkungen besprochen werden k�nnen.\\
\\
{\bf Wichtig:} Die Vortr�ge sind vorab mit den Dozierenden zu besprechen. Buchen Sie dazu bis sp�testens eine Woche vor Ihrem Termin einen Zeitslot �ber \href{https://cal.meetergo.com/joaquin-veith/30-min-meeting-or-joaquin-veith}{diesen Link} (nur Joaquin Veith). Wir werden dann gemeinsam in Pr�senz (in 1R18 bei Joaquin Veith und 1R22 f�r Fabian Hennig) Ihren Vortrag besprechen. Schicken Sie dazu Ihren Foliensatz und ihr Handout vorab via Mail.

\end{document}
